\documentclass[pdflatex,sn-mathphys-num]{sn-jnl}

\usepackage{graphicx}%
\usepackage{multirow}%
\usepackage{amsmath,amssymb,amsfonts}%
\usepackage{amsthm}%
\usepackage{mathrsfs}%
\usepackage[title]{appendix}%
\usepackage{xcolor}%
\usepackage{textcomp}%
\usepackage{manyfoot}%
\usepackage{booktabs}%
\usepackage{algorithm}%
\usepackage{algorithmicx}%
\usepackage{algpseudocode}%
\usepackage{listings}%
\usepackage{lmodern}%
\usepackage{siunitx}
\usepackage{microtype}
\usepackage{derivative}
\usepackage[version=4]{mhchem}
\usepackage{cleveref}
\usepackage[spanish]{babel}

\renewcommand{\arraystretch}{1.2}
\sisetup{group-digits=true,
  group-separator={\,},
  separate-uncertainty
}

\theoremstyle{thmstyleone}%
\newtheorem{theorem}{Teorema}%
\newtheorem{proposition}[theorem]{Proposición}%

\theoremstyle{thmstyletwo}%
\newtheorem{example}{Ejemplo}%
\newtheorem{remark}{Observación}%

\theoremstyle{thmstylethree}%
\newtheorem{definition}{Definición}%

\raggedbottom
\unnumbered% uncomment this for unnumbered level heads

\begin{document}

\title[Análisis de Brechas en Semiconductores]{Análisis Fundamental de las
Brechas de Energía en Semiconductores y el Efecto del Desorden Topológico}

\author{\fnm{Julian L.} \sur{Avila-Martinez}}
\email{jlavilam@udistrital.edu.co}

\affil{\orgdiv{Programa Académico de Física}, \orgname{Universidad Distrital
Francisco José de Caldas}, \orgaddress{\city{Bogotá D.C.},
\country{Colombia}}}

\maketitle

\section{Transiciones Energéticas y Brechas Prohibidas}
\label{sec:brechas}

La distinción fundamental entre un semiconductor de brecha ancha y uno de brecha
estrecha reside en la topología del espectro energético y en la restricción
impuesta sobre las amplitudes de transición interbanda. En un sistema físico
donde la separación espectral excede los \qty{4}{\eV}, el comportamiento se
aproxima de manera estricta al de un dieléctrico o aislante.
Materiales semiconductores con una brecha energética considerable, como el
dióxido de titanio cuyo valor oscila entre \qty{3.0}{\eV} y \qty{3.3}{\eV},
presentan una distribución térmica donde la probabilidad de ocupación de estados
en la banda de conducción es prácticamente nula frente a fluctuaciones térmicas
habituales.
Consecuentemente, se requiere una interacción con fotones de alta energía,
correspondientes al espectro ultravioleta, para inducir transiciones. Esto le
confiere al material una alta transmitancia de la radiación en el régimen del
espectro electromagnético visible.

En el caso de los semiconductores intrínsecos de brecha estrecha, tales como el
silicio y el germanio, existen estados localizados adyacentes al nivel de Fermi
que mitigan el requerimiento energético para la promoción de un electrón desde
la banda de valencia. La reducción de la escala energética
facilita el incremento sustancial en la densidad de estados operacionales, lo
que dota al material de una mayor susceptibilidad a factores termodinámicos y
electromagnéticos ambientales sin precisar estimulaciones de gran magnitud.

\section{Ruptura de la Estructura Topológica en Sistemas Amorfos}
\label{sec:amorfos}

La complejidad intrínseca en la descripción del espectro energético de
los materiales amorfos es una consecuencia directa de la pérdida de invarianza
traslacional global. En el límite de un sólido cristalino ideal, el operador
hamiltoniano es invariante ante el grupo algebraico de traslaciones espaciales
discretas. Esta simetría permite estructurar las soluciones del sistema de
conformidad con el teorema de Bloch. Las funciones de estado
constituyen representaciones irreducibles, operando una transformación canónica
desde la red geométrica en el espacio real hacia la red recíproca en el espacio
de momentos. En este formalismo, las estructuras de banda
se obtienen unívocamente evaluando el espectro a lo largo de caminos de alta
simetría circunscritos en la zona de Brillouin.

En contraposición, una red amorfa exhibe un desorden que trunca la periodicidad
espacial macroscópica; transcurridas pocas magnitudes a partir de una celda
unitaria de referencia, la organización molecular diverge. La
anulación del grupo de simetría de traslación global implica que el momento
cristalino $k$ carece de conservación, y por extensión, la noción topológica
de la zona de Brillouin deja de existir. Al invalidarse las
condiciones del teorema de Bloch, los estados propios no pueden describirse por
medio de ondas planas moduladas, exigiendo un tratamiento restringido a un orden
de corto alcance o cristalografía local. Debido a esto, el
análisis requiere el uso de formulaciones fundamentadas en la densidad
estadística, como el modelo de Tauc para redes amorfas, o los
modelos de transporte fenomenológicos de Mott que examinan la localización de
portadores de carga en presencia del desorden local.

\end{document}
